\documentclass[11pt]{article}
\usepackage[margin=1in]{geometry}
\usepackage{amsmath,amssymb,amsthm}
\usepackage{hyperref}
\newtheorem{theorem}{Theorem}
\newtheorem{lemma}[theorem]{Lemma}
\newtheorem{proposition}[theorem]{Proposition}
\theoremstyle{remark}
\newtheorem{remark}[theorem]{Remark}

\title{Disproof of Conjecture 00000002141:\\
there are two different cospectral pairs of 4-regular graphs of order 10}
\author{jilint777}
\date{2026-10-04}

\begin{document}
\sloppy
\maketitle

\begin{abstract}
Conjecture 00000002141 asserts that the minimal order of a pair of cospectral
(non-isomorphic) regular graphs is $10$, and that the cospectral regular pair of
order $10$ is unique. We exhibit two cospectral pairs $(A_1,B_1)$ and
$(A_2,B_2)$ of 4-regular graphs on $10$ vertices. The two pairs have
different spectra, and no graph of the first pair is isomorphic to a graph
of the second. So the uniqueness clause is false. An exhaustive enumeration
shows that order $10$ has exactly two 4-regular cospectral pairs and their two
5-regular complements. The disproof is formalized in Lean~4 (core library
only).
\end{abstract}

\section{The conjecture}
\begin{quote}
\emph{Definition:} Cospectral regular graphs are non-isomorphic regular graphs
with the same adjacency spectrum. \emph{Conjecture:} The minimal order of a
cospectral regular pair is 10 (the Saltire pair); the cospectral regular pair
of order 10 is unique (classification uniqueness).
\end{quote}
The Chinese version says the same (the order-10 pair is the only one). The conjecture
is a conjunction, so it suffices to refute the uniqueness clause. We read
``the pair is unique'' in the weakest sense: any two cospectral regular pairs
$\{G,H\}$ and $\{G',H'\}$ of order $10$ coincide up to isomorphism as unordered
pairs. (The first clause is in fact true; see Section~\ref{sec:enum}. The name
``Saltire pair'' usually refers to the non-regular pair $K_{1,4}$,
$C_4\cup K_1$ on 5 vertices, but this plays no role.)

\section{The graphs}
All four graphs have vertex set $\{0,\dots,9\}$ and $20$ edges.
\begin{align*}
A_1:&\ 03,05,07,09,14,15,17,19,25,26,27,28,36,38,39,46,48,49,57,68\\
B_1:&\ 03,05,07,08,14,15,18,19,25,26,27,29,36,37,38,46,48,49,57,69\\
A_2:&\ 03,05,07,08,14,15,17,19,25,26,28,29,36,37,38,46,48,49,57,69\\
B_2:&\ 03,04,06,07,15,16,18,19,25,27,28,29,34,36,37,48,49,57,58,69
\end{align*}
(In graph6 format: \verb|ICR`vGyu?|, \verb|ICR`tjWY_|, \verb|ICR`uiwY_|,
\verb|ICdedhkY_|.) Each vertex has degree $4$ in each graph.

\section{Cospectrality}
For a graph with adjacency matrix $A$ and eigenvalues
$\lambda_1,\dots,\lambda_n$, the number of closed walks of length $k$ is
$\operatorname{tr}(A^k)=\sum_i\lambda_i^k=p_k$. By Newton's identities
$ke_k=\sum_{i=1}^k(-1)^{i-1}e_{k-i}p_i$, the power sums $p_1,\dots,p_n$
determine the elementary symmetric functions $e_1,\dots,e_n$, hence the
characteristic polynomial $\prod(x-\lambda_i)$, hence the spectrum. The
converse is clear. So two graphs on $n$ vertices are cospectral if and only if
they have equally many closed walks of each length $k\le n$.

\begin{proposition}
$\operatorname{tr}(A_1^k)=\operatorname{tr}(B_1^k)$ and
$\operatorname{tr}(A_2^k)=\operatorname{tr}(B_2^k)$ for $0\le k\le10$. Equivalently,
\begin{align*}
\det(xI-A_1)=\det(xI-B_1)={}&x^{10}-20x^8-16x^7+110x^6+136x^5\\&-180x^4-320x^3+9x^2+200x+80,\\
\det(xI-A_2)=\det(xI-B_2)={}&x^{10}-20x^8-14x^7+108x^6+104x^5\\&-183x^4-188x^3+80x^2+68x-16.
\end{align*}
The two polynomials differ. For example, $\operatorname{tr}(A^3)$ is six times the
number of triangles, and $\operatorname{tr}(A_1^3)=48$ ($8$ triangles) while
$\operatorname{tr}(A_2^3)=42$ ($7$ triangles).
\end{proposition}
\begin{proof}
Direct computation, in Lean (\texttt{cospectral\_1}, \texttt{cospectral\_2},
\texttt{spectra\_differ}) and in \texttt{verify.py}. The script computes the
characteristic polynomials in two independent exact ways: Faddeev--LeVerrier,
and $\det(xI-A)$ by exact Gaussian elimination at $x=0,\dots,10$.
\end{proof}

\section{Non-isomorphism}
Call an edge $uv$ \emph{triangle-free} if $u$ and $v$ have no common
neighbour. Both of the following properties are preserved by graph
isomorphisms:
\begin{itemize}
\item[(P)] some edge is triangle-free;
\item[(Q)] some triangle-free edge shares no endpoint with any other
triangle-free edge.
\end{itemize}
\begin{lemma}
$A_1$ fails (P), while $B_1$, $A_2$ and $B_2$ satisfy (P). $A_2$ satisfies (Q),
while $B_2$ fails (Q).
\end{lemma}
\begin{proof}
The triangle-free edges are: none in $A_1$; $15,36$ in $B_1$; $25,28,36,48$ in
$A_2$, where $36$ is isolated; and $29,48,49$ in $B_2$, which form the path
$2\!-\!9\!-\!4\!-\!8$.
\end{proof}
Hence $A_1\not\cong B_1$ and $A_2\not\cong B_2$, so both are cospectral regular
pairs. Also $A_1\not\cong A_2$ and $A_1\not\cong B_2$, so the unordered pairs
$\{A_1,B_1\}$ and $\{A_2,B_2\}$ are different up to isomorphism. The two pairs
also have different spectra.

\begin{theorem}
The cospectral regular pair of order $10$ is not unique. Conjecture 00000002141
is false.
\end{theorem}

\begin{remark}
The two pairs are not complements of each other: all four graphs are
4-regular, while complements of 4-regular graphs on 10 vertices are 5-regular.
So uniqueness fails even up to complementation.
\end{remark}

\section{Enumeration}\label{sec:enum}
With nauty's \texttt{geng}, \texttt{verify.py} enumerates all $k$-regular graphs
on $n\le10$ vertices and groups them by characteristic polynomial. There is
no cospectral pair of regular graphs on fewer than $10$ vertices. On $10$
vertices there are exactly two cospectral pairs of 4-regular graphs, namely the
two above (out of $60$ 4-regular graphs), and exactly two of 5-regular graphs,
namely their complements. So the first clause (minimal order $10$) is true, and
there are exactly four cospectral regular pairs of order $10$.

\section{Lean formalization}
The directory \texttt{lean4/} is a self-contained Lean~4.19.0 project (core
library only, no Mathlib).
\begin{itemize}
\item Graphs on $\{0,\dots,9\}$ are Boolean adjacency functions.
\texttt{IsSimple} means symmetric and loopless, and \texttt{Regular g k} means
every vertex has degree $k$.
\item \texttt{closedWalks g k} is $\operatorname{tr}(A^k)$, computed with explicit
matrix products. \texttt{Cospectral g h} means equal closed-walk counts for
all $k\le10$, which by the argument above is equivalent to having the same
adjacency spectrum.
\item \texttt{Iso g h} is the existence of a bijection $\sigma$ of
$\{0,\dots,9\}$ (with inverse $\tau$) such that $g(u,v)=h(\sigma u,\sigma v)$.
\texttt{iso\_hasNT} and \texttt{iso\_hasIsolatedNT} prove that (P) and (Q) are
isomorphism invariants. \texttt{invariants} evaluates (P) and (Q) on the four
graphs.
\item \texttt{CospectralRegularPair g h} combines simple, regular, cospectral
and non-isomorphic. \texttt{UniquePair} says any two such pairs coincide up to
isomorphism, in either order. \texttt{conjecture\_00000002141\_false} proves
$\neg$\texttt{UniquePair}, using \texttt{pair\_1}, \texttt{pair\_2},
\texttt{not\_iso\_A1\_A2} and \texttt{not\_iso\_A1\_B2}.
\end{itemize}
\texttt{lake build} succeeds. There is no \texttt{sorry}, no
\texttt{native\_decide}, and no added axioms. \texttt{\#print axioms} shows at
most \texttt{propext} and \texttt{Quot.sound}.

\section{Reproduction}
\begin{verbatim}
cd lean4 && lake build     # Lean 4.19.0, core only
python3 verify.py          # stdlib (+ nauty-geng if present)
pdflatex report.tex
\end{verbatim}

\begin{thebibliography}{9}
\bibitem{bh} A.~E.~Brouwer and W.~H.~Haemers, \emph{Spectra of Graphs},
Springer, 2012 (cospectrality and closed walks; GM switching).
\bibitem{hs} W.~H.~Haemers and E.~Spence, Enumeration of cospectral graphs,
\emph{European J. Combin.} 25 (2004), 199--211.
\end{thebibliography}
\end{document}
